\chapter*{How to use this guide}
\addcontentsline{toc}{chapter}{How to use this guide}

This guide explains, in more detail than a paper would, what the forward model of Chapter~3 does, why each modelling decision was taken, what the alternatives were, what evidence supports the choice, and where in the code each piece lives.
It is written to be read alongside the chapter, not instead of it: the chapter states the model, and this guide is the reasoning behind it.

Every chapter follows the same pattern.
\begin{enumerate}[nosep]
\item The theory, derived step by step, with every force balance written as \emph{force up = force down}.
\item The decisions, each in a green box: what we chose, what else we could have done, and why.
\item The defence, in an orange box: the questions an examiner is likely to ask, with the answer and the number that supports it.
\item The code, in a blue box: which file and which function implement the step, and what they compute.
\end{enumerate}

Chapter~\ref{ch:ledger} collects every decision in one table, with its evidence and its caveat, and ends with the questions most likely to come up in a viva.
If you read only one chapter before a meeting, read that one.

\section*{The model in one page}

We want the density of the mantle.
We observe two things at the surface: the height of the seafloor (and, later, of the land), and the geoid.
A density anomaly in the mantle makes the mantle flow; the flow pushes or pulls on the base of the lithosphere, which deflects the surface; and every mass that is displaced, the anomaly itself and the boundaries it deflects, changes the gravity field and so the geoid.
All of this is linear in the density anomaly, so the forward problem is a linear operator, and with Gaussian priors and errors the inverse problem has an exact, analytic posterior.

We build the forward model up in stages, making it more physical one assumption at a time; each stage can be checked against the one before by a forward comparison.
\begin{enumerate}[nosep]
\item \textbf{The core model} (Chapter~\ref{ch:core}): mantle flow below a lithosphere with no strength, covered everywhere by water that follows the geoid.
It predicts oceanic residual depths and the geoid, and it is what the first inversion uses.
\item \textbf{An elastic plate on land and sea} (Chapter~\ref{ch:plate}): the lithosphere bends, and the continents are dry.
This is what lets continental data into the final inversion.
\item \textbf{The lithosphere as an unknown} (Chapters~\ref{ch:lithosphere} and~\ref{ch:geoid}): the crust and lithospheric mantle become unknowns about a reference model, WINTERC-G, and the data become total topography.
\item \textbf{The real ocean and the sea-level constant} (Chapter~\ref{ch:ocean}): water only where the oceans are, solved together with the deflection.
\end{enumerate}
Each stage reduces to the previous one in a stated limit, and the difference is measured, so each added piece of physics is justified by a number rather than by assertion.
There are two inversions: the core model with oceanic residual depths and the geoid, and the full model with total topography and the geoid.
The stages in between build the forward model; they are not inverted separately.

\section*{What has been verified}
Every equation in the forward model has been checked against something independent of the code that implements it: a closed-form solution, a symbolic derivation, a second solver written separately (\texttt{fem\_1d}), a brute-force numerical integral, or an exact conservation law.
Chapter~\ref{ch:assembly} lists the checks and their results; the smallest disagreement is at the level of rounding error and the largest, for the geoid of a compensated root computed two entirely different ways, is \qty{3}{\percent}, a gap not yet explained.
An independent audit of the physics and the code, and three simulated referee reports, found errors in the code, all fixed and described where they occur; a later independent review of this guide found and corrected errors in its explanations.
